\section{Conclusion}
\label{sec:conclusion}

We compared post-hoc, training-time, and stacked embedding compression for dense
retrieval on common ground, across three encoders with different retrieval
priors. No single method dominates. The right choice depends on whether the
encoder is pre-trained for retrieval and on the memory budget: post-hoc PQ leads
for the retrieval-pre-trained encoder, while training-time binarization leads for
the general-purpose encoders, and stacked MRL$+$PQ leads at high compression for
all. The most surprising result is that binarization-aware training can exceed
the full-precision baseline on an encoder without retrieval pre-training. We
explained this through effective dimension: binarization spreads embedding
variance across more axes, improving ranking without changing recall, and the
effect is structural rather than a symptom of under-training.

Future work follows directly. Scaling the evaluation to full BEIR and MTEB, and
adding measured latency and index size, would test whether the backbone-dependent
frontiers hold at deployment scale. Pairing MRL with recurrent residual
binarization~\citep{gan2023binary} would give runtime-selectable dimension and
bit depth. And a controlled anisotropy intervention would turn the
effective-dimension correlation into a causal account of when compression helps
retrieval.
